\subsection{The exact unbordered range on each stable-graph stratum}
\label{sec:A1-range}

This subsection concerns the unbordered construction: the external
marked points remain cusps. Let $\Gamma$ be a stable graph with edge
set $E$, and fix its marked normalization $(Y_v)_{v\in V}$.
At a vertex, $n_v$ counts all node branches and external markings;
assume $2g_v-2+n_v>0$. Write $I_v$ for its node branches, and let
\[
 A_v:\R^E\longrightarrow\R^{I_v},\qquad
 (A_vs)_a=s_{e(a)}.
\]
A loop edge has two distinct branches at the same vertex and therefore
occurs twice among these coordinates. For each edge the two geodesic
boundary lengths are required to be equal, with common value
$2\pi\sqrt{s_e}$.

Let $S_v$ and $K_v$ be the potential and support function of
Theorems~\ref{A2b-energy}--\ref{A2b-image}, applied to $Y_v$ with
active set $I_v$ and the other markings as cusps. A vertex with no
active branches contributes zero. Define
\begin{equation}\label{A1-graph-potential}
 \Psi_\Gamma(s)=\sum_v S_v(A_vs),\qquad
 K_\Gamma(p)=\sum_v K_v(A_vp).
\end{equation}
All logarithmic norms below use arbitrary fixed tangent frames at the
branches; changing the frames translates their image in $\R^E$.

\begin{theorem}[Proved: exact range, tangent map, and fundamental group]
\label{A1-graph-range}
For the fixed normalization, the unbordered canonical surgery map,
with actual toric seam identifications as its phase variables, is a
real-analytic diffeomorphism onto
\[
 \mathcal U_\Gamma(Y)=
 \left\{(q_e)\in\prod_{e\in E}(T_{a_e}Y\otimes T_{b_e}Y)^\times:
   \sum_e p_e\log|q_e|<K_\Gamma(p)
   \quad\text{for all }0\ne p\in[0,\infty)^E\right\}.
\]
Its logarithmic radial image is open, convex, and downward closed.
If $E\ne\varnothing$, this is a proper subset of the unrestricted
nonzero-tensor fiber. The tangent map of surgery is an isomorphism.

As the normalization varies on its fixed graph stratum, these images
form an open real-analytic subbundle of the nonzero-tensor bundle.
Its inclusion is an equivariant homotopy equivalence, and consequently
induces an isomorphism on fundamental groups, including orbifold
fundamental groups when finite automorphisms are retained. These
conclusions allow loop edges as well as edges between different
components.
\end{theorem}
\begin{proof}
At the branch $a$ of an edge $e$, let $r_a(s)$ denote the canonical
cap radius obtained on $Y_{v(a)}$ with parameters $A_{v(a)}s$.
The tensor computation in Lemma~\ref{A1-tensor} gives
\[
 \log|q_e|=\log r_{a_e}(s)+\log r_{b_e}(s)
          =\frac{\partial\Psi_\Gamma}{\partial s_e}(s).
\]
This chain rule counts both branches of a loop separately. Its Hessian
is
\[
 d^2\Psi_\Gamma=\sum_v A_v^T(d^2S_v)A_v.
\]
Each nonempty term is positive definite on the branch variables. If
$\dot s\ne0$, choose an edge with $\dot s_e\ne0$; at least one
$A_v\dot s$ is nonzero. Hence the displayed Hessian is positive
definite on $\R^E$. The logarithmic tensor map is therefore locally
invertible and globally injective by integration along straight
segments, exactly as in Theorem~\ref{A2b-energy}.

For completeness the same argument gives the entire stated image,
not just an injective parametrization. Its recession function is
$K_\Gamma$: apply Lemma~\ref{A2b-recession} to each $A_vs_0+tA_vp$
and sum. The supports are finite and positively homogeneous, with
$K_v(0)=0$. Every actual radius vector satisfies all the strict
inequalities, because at least one nonzero $A_vp$ gives a strict
component inequality.

Conversely, let $x\in\R^E$ satisfy those strict inequalities.
The convex function $\Psi_\Gamma(s)-x\cdot s$ has bounded sublevels:
an escaping sequence, normalized in direction, would give a nonzero
$p\geq0$ with recession slope $K_\Gamma(p)-p\cdot x\leq0$, by
the segment argument in the proof of Theorem~\ref{A2b-image}.
Its lower semicontinuous extension to the closed orthant attains its
minimum. A finite boundary minimum is impossible. Along an inward
diagonal, every vanishing edge variable produces two cap terms tending
to $-\infty$ by the uniform collar bound, while all other cap terms
have uniform finite upper bounds. The inward derivative thus tends
to $-\infty$. The minimum is interior, unique, and has gradient $x$.
This proves the asserted logarithmic range and its analytic inverse.

Normalize the cap maps to have positive derivatives in the chosen
tangent frames. The gluing identification is then
$w_{a_e}w_{b_e}=c_e$, $|c_e|=1$, and the argument of $q_e$ is
precisely that of $c_e$. Thus all phases occur independently and
uniquely. Equal lengths make this identification an isometry, so
Proposition~\ref{A1-retained-discs} gives the unique welded hyperbolic
surface with its indicated cut. This proves the real-analytic
diffeomorphism. If the normalization also varies, the map preserves
the normalization coordinate. Its derivative is block triangular
with identity on those coordinates, the positive Hessian on length
variables, and identity on phase variables. It remains invertible;
the local inverses fit together by the fiberwise uniqueness.

To see properness when there is an edge, take $p$ to be its coordinate
unit vector. The corresponding inequality bounds $\log|q_e|$ above
by the finite number $K_\Gamma(p)$. The unrestricted tensor fiber
has no such bound.

Finally set $s_e=1$ for every edge. The canonical caps of these common
lengths determine an intrinsic circle in each nodal tensor line, by
taking the product of their shadow boundary circles. Let $\mathcal C$
be the product of these circles over the normalization stratum.
For any nonzero tensor tuple, let $a_e>0$ be its radial ratio to the
corresponding circle. The homotopy
\[
 H_t((q_e))=(a_e^{-t}q_e)_{e\in E},\qquad0\leq t\leq1,
\]
is a strong deformation retraction of the full nonzero-tensor bundle
onto $\mathcal C$. It also retracts $\mathcal U_\Gamma$ onto
$\mathcal C$, because in logarithmic coordinates it follows a segment
to an interior point of a convex fiber. The construction is intrinsic
and respects graph automorphisms, markings, and relabellings. Both
retractions descend to the relevant orbifold quotients. Their common
retract proves the homotopy and fundamental-group assertions.
For a fixed normalization this reduces to the identity on the
$\Z^E$ fundamental group of the phase torus.
\end{proof}

\begin{corollary}[Proved: the two smallest one-node ranges]
\label{A1-small-ranges}
For the separating one-node $(0,4)$ curve, normalize each rational
component by putting its node at $0$ and its two external markings at
$1,\infty$. In the resulting tensor frame the exact range is
\[
                         0<|q|<16.
\]
For the irreducible one-node $(1,1)$ curve, put its two branches at
$0,1$ and its marking at $\infty$ on its rational normalization.
In the frame $\partial_z|_0\otimes\partial_z|_1$ the exact range is
\[
                         0<|q|<1.
\]
In both cases every phase occurs and the length is uniquely determined
by $|q|$.
\end{corollary}
\begin{proof}
For $(0,4)$ each of the two components has one active cap. Its support
value at weight $1$ is $\log4$ by the Koebe limit in
Corollary~\ref{A2-explicit-limits}. Thus
$K_\Gamma(p)=p\log16$, and Theorem~\ref{A1-graph-range} gives the
first range. Equivalently the two identical one-cap radii increase
from $0$ to $4$ with their common length.

For $(1,1)$ there is one component and a loop edge. Its incidence map
is $s\mapsto(s,s)$. The two-weight Strebel discs at $0,1$ are the
half-planes separated by $\operatorname{Re}z=1/2$, each of conformal
radius $1$. Hence $K_v(p,p)=0$ and the logarithmic range is
$\log|q|<0$. This proves the second assertion. Strict monotonicity
is the one-dimensional positive Hessian of $S_v(s,s)$, and the
phase statement is the seam tensor formula.
\end{proof}

\paragraph{Marking convention and bordered scope.}
The phase in this statement is an actual orientation-reversing toric
identification of the two boundary circles, so it belongs to $S^1$.
If a fully marked Fenchel--Nielsen twist is instead kept as a real
number, its full Dehn-twist translates give the same tensor phase;
one must take that quotient, or lift the tensor phase to the matching
covering space, before asserting the displayed diffeomorphism.

For the bordered target in the manuscript, the external admissible
disc domains are retained as data, not merely their boundary lengths.
Fixing their lengths while varying internal node lengths generally
changes those external domains, by the cap-dependence result in
A.2(a). Therefore Theorem~\ref{A1-graph-range} does not by itself
identify the image in a fiber where the full external domains have
been fixed. Proposition~\ref{A1-bordered-landing} proves the further
obstruction that external-only admissibility need not survive
discarding the node caps. Theorem~\ref{A1-bordered-shadow-range}
instead includes external shadows among the joint gradient variables,
and Corollary~\ref{A1-bordered-shape-image} gives the resulting exact
criterion when the external domains themselves are prescribed.
